\subsection{任意超越扩张的情形}

在本号中，我们使用下述记号：K 是域，\(K'\) 是 K 的扩域，v 是 K 上的赋值，\(v'\) 是 v 到 \(K'\) 的延拓，\(\Gamma\) 和 k（分别地，\(\Gamma'\) 和 \(k'\)）分别是 \(v\)（分别地，\(v'\)）的阶群和剩余域。我们记：

\[
d \left(\mathrm{K} ^ {\prime} / \mathrm{K}\right) = \dim . \mathrm{al} _ {\mathrm{K}} \mathrm{K} ^ {\prime} = \text { transcendence   degree   of } \mathrm{K} ^ {\prime} \text { over } \mathrm{K};
\]

\[
s (v ^ {\prime} / v) = \dim . \mathrm{al} _ {k} k ^ {\prime} = \text { transcendence   degree   of } k ^ {\prime} \text { over } k;
\]

\[
r (v ^ {\prime} / v) = r (\Gamma^ {\prime} / \Gamma) = \text { rational   rank   of } \Gamma^ {\prime} / \Gamma ,
\]

若等号右边各项有限；否则约定 \(d(\mathbf{K}' / \mathbf{K}) = +\infty\)（分别地，\(s(v' / v) = +\infty\)、\(r(v' / v) = +\infty\)）。

定理 1. 设 \(x_1, \ldots, x_s\) 是 \(v'\) 的环中的元素，其典范像 \(\bar{x}_i\) 在 \(k'\) 中在 \(k\) 上代数独立；设 \(y_1, \ldots, y_r\) 是 \(K'\) 中的元素，使得 \(v'(y_j)\) 在 \(\Gamma'/\Gamma\) 中的典范像在 \(Z\) 上线性无关。则 \(r + s\) 个元素 \(x_1, \ldots, x_s, y_1, \ldots, y_r\) 属于 \(K'\)，并在 \(K\) 上代数独立；\(v'\) 限制于 \(K(x_1, \ldots, x_s, y_1, \ldots, y_r)\) 后的剩余域为 \(k(\bar{x}_1, \ldots, \bar{x}_s)\)，并且

\[
\Gamma + \mathbf {Z} v ^ {\prime} (y _ {1}) + \dots + \mathbf {Z} v ^ {\prime} (y _ {r})
\]

作为有序群。

当 \(r + s = 0\) 时断言显然。我们对 \(r + s\) 作归纳。若 \(r' \leqslant r\)、\(s' \leqslant s\) 且 \(r' + s' < r + s\)，则归纳假设表明：当 \(K\) 换成 \(K(x_1, \ldots, x_s, y_1, \ldots, y_{r'})\)，且族 \((x_1, \ldots, x_s), (y_1, \ldots, y_r)\) 换成 \((x_{s'+1}, \ldots, x_s), (y_{r'+1}, \ldots, y_r)\) 时，定理 1 的假设仍成立。因此问题归结为下述两种情形之一：

\begin{enumerate}
\def\labelenumi{(\alph{enumi})}
\item
  存在 \(x\) 属于 \(v'\) 的环，使得 \(\bar{x}\) 在 \(k\) 上超越；则须证明 \(x\) 在 \(\mathbf{K}\) 上超越，且 \(v'\) 限制于 \(\mathbf{K}(x)\) 后的剩余域为 \(k(\bar{x})\)，阶群为 \(\Gamma\)。
\item
  存在 \(y\) 属于 \(\mathbf{K}'\)，使关系 \(nv'(y) \in \Gamma\) 和 \(n \in \mathbf{Z}\) 蕴含 \(n = 0\)；则须证明 \(y\) 在 \(\mathbf{K}\) 上超越，且 \(v'\) 限制于 \(\mathbf{K}(y)\) 后的剩余域为 \(k\)，阶群为 \(\Gamma + \mathbf{Z}v'(y)\)。
\end{enumerate}

现在第 8 节第 1 节的命题 1 表明，\(x\)（分别地，\(y\)）不可能在 \(K\) 上代数。(a)（分别地，(b)）的其余断言由第 1 节的命题 2（分别地，命题 1）立即得到。

推论 1. 不等式

\[
s (v ^ {\prime} / v) + r (v ^ {\prime} / v) \leqslant d (\mathbf {K} ^ {\prime} / \mathbf {K})\tag{9}
\]

成立。

此外，若 \(K'\) 是 K 的有限生成扩张，且 (9) 中取等号，则 \(\Gamma'/\Gamma\) 是有限生成 Z-模，且 \(k'\) 是 k 的有限生成扩张。

令 \(r\) 和 \(s\) 为自然数，满足 \(r \leqslant r(v'/v)\) 且 \(s \leqslant s(v'/v)\)；我们证明 \(r + s \leqslant d(\mathbf{K}'/\mathbf{K})\)，这将证明 (9)。根据假设，存在 \(x_1, \ldots, x_s, y_1, \ldots, y_r\) 属于 \(\mathbf{K}'\) 的元素，它们满足定理 1 的假设。

所以它们在 K 上代数独立，这就表明不等式 \(r + s \leqslant d(\mathbf{K}'/\mathbf{K})\) 成立。

若 \(\mathbf{K}'\) 是 \(\mathbf{K}\) 的有限生成扩张，\(d(\mathbf{K}' / \mathbf{K})\) 有限，因而 \(s(v' / v)\) 和 \(r(v' / v)\) 也有限；我们分别记为 \(s\) 和 \(r\)。存在元素 \(x_{1},\ldots ,x_{s},\) \(y_{1},\ldots ,y_{r}\) 属于 \(\mathbf{K}'\)，它们满足定理 1 的假设。若 \(r + s = d(\mathbf{K}' / \mathbf{K})\)，这些元素构成 \(\mathbf{K}'\) 在 \(\mathbf{K}\) 上的超越基，因而 \(\mathbf{K}'\) 是 \(\mathbf{K}'' = \mathbf{K}(x_1,\dots,y_r)\) 的有限代数扩张。令 \(\Gamma''\) 和 \(k''\) 为 \(v'\) 限制于 \(\mathbf{K}''\) 所得赋值的阶群和剩余域。根据定理 1，\(\Gamma'' / \Gamma\) 是有限生成的 \(\mathbf{Z}\)-模，且 \(k''\) 是 \(k\) 的有限生成纯扩张。另一方面，由于 \(\mathbf{K}'\) 是 \(\mathbf{K}''\) 的有限代数扩张，\(\Gamma' / \Gamma\) 是有限群，且 \(k'\) 是 \(k''\) 的有限代数扩张（§ 8，第 1 节，引理 2）。这证明了推论。

推论 2. 令 \(h\) 和 \(h'\) 分别为 \(v\) 和 \(v'\) 的高。则

\[
s (v ^ {\prime} / v) + h ^ {\prime} \leqslant d (\mathrm{K} ^ {\prime} / \mathrm{K}) + h.\tag{10}
\]

由命题 3，\(h' \leqslant r(v'/v) + h\)。

推论 3. 设 \(\mathbf{K}'\) 是 \(\mathbf{K}\) 的有限生成扩张，\(\Gamma\) 同构于按字典序排列的 \(\mathbf{Z}^h\)，且公式 (10) 中取等号。则 \(\Gamma'\) 同构于按字典序排列的 \(\mathbf{Z}^{h'}\)，且 \(k'\) 是 \(k\) 的有限生成扩张。

若 (10) 中取等号，则 (9) 中也取等号，从而 \(k'\) 是 k 的有限生成扩张，且 \(\Gamma'\) 是有限生成 Z-模。此外，比较 (9) 与 (10)，可见 \(h' - h = r(\Gamma'/\Gamma)\)，从而 \(h' = r(\Gamma')\)，于是第 2 节的命题 4 表明 \(\Gamma'\) 同构于按字典序排列的 \(Z^{h'}\)。

推论 4. 设 \(v\) 是非真赋值（此时 \(k = K\)）。则

\[
h (\Gamma^ {\prime}) + d (k ^ {\prime} / \mathbf {K}) \leqslant r (\Gamma^ {\prime}) + \mathbf {D} (k ^ {\prime} / \mathbf {K}) \leqslant d (\mathbf {K} ^ {\prime} / \mathbf {K}).\tag{11}
\]

特别地，若 \(v'\) 的高为 1，则

\[
d (k ^ {\prime} / \mathbf {K}) \leqslant d (\mathbf {K} ^ {\prime} / \mathbf {K}) - 1;\tag{12}
\]

此外，若 \(\mathbf{K}'\) 是 \(\mathbf{K}\) 的有限生成扩张，且 (12) 中取等号，则 \(v'\) 是离散赋值，且 \(k'\) 是 \(\mathbf{K}\) 的有限生成扩张。

这是推论 1、2、3 的一系列特殊情形。


